% !PRINTED 637
\weberpage{637}
%<<CONT
where $m = n'n''$, and the $x_{\nu,\nu'}$ are integral rational functions of
$z$. From (7) there results, by elimination of the $\omega_i$:

\[
\begin{vmatrix}
x_{1,1} - \omega, & x_{1,2}, & \ldots, & x_{1,m} \\
x_{2,1}, & x_{2,2} - \omega, & \ldots, & x_{2,m} \\
\hdotsfor{4} \\
x_{m,1}, & x_{m,2}, & \ldots, & x_{m,m} - \omega
\end{vmatrix} = 0,
\]

and this is an equation for $\omega$ of the form (5).

By repeated application of this theorem it follows that every
\gesperrt{integral rational function of integral functions is again
an integral function.}

7. \gesperrt{An integral function $\omega$ is called divisible by an
integral function $\omega'$ if there exists a third integral function
$\omega''$ such that}

\weq{8}{\omega = \omega'\omega''}

\gesperrt{holds.}

From this definition it follows at once:

\gesperrt{If $\omega$ is divisible by $\omega'$ and $\omega'$ is divisible by $\omega''$, then
$\omega$ is also divisible by $\omega''$.}

\gesperrt{If $\omega'$ and $\omega''$ are divisible by $\omega$, then $\omega' \pm \omega''$
is also divisible by $\omega$.}

\gesperrt{If $\omega_1, \omega_2, \omega_3, \ldots$ are divisible by $\omega$ and $\omega_1', \omega_2', \omega_3', \ldots$
are arbitrary integral functions, then}

\[
\omega_1 \omega_1' + \omega_2 \omega_2' + \omega_3 \omega_3' + \cdots
\]

\gesperrt{is also divisible by $\omega$.}

\begin{center}
\textbf{§ 176. Minimal basis and field discriminant.}
\end{center}

Since every function of the field $\Omega$ can be converted by multiplication
with an integral rational function into an integral algebraic
function of $z$, there also exist field bases
consisting of \gesperrt{integral functions} (e.g., in the notation
 of § 171, (2), the powers of $a_0 \theta$).

If now

\weq{1}{\omega_1, \ \omega_2, \ \ldots, \ \omega_n}

is such a basis consisting of integral functions, then every
in the form

\weq{2}{\omega = x_1 \omega_1 + x_2 \omega_2 + \cdots + x_n \omega_n}
%CONT>>
